\documentclass[11pt]{article}
\usepackage{fontspec}
\usepackage[french,english]{babel}
\usepackage{amsmath}
\usepackage{amssymb}
\usepackage{amsthm}
\usepackage{mathrsfs}
\usepackage{tikz-cd}
\usepackage{enumitem}
\usepackage{lmodern}

\newcommand{\printedpage}[1]{\marginpar{\footnotesize p.~#1}}
\newcommand{\msp}[1]{\marginpar{\footnotesize [#1]}}
\newcommand{\HOTlc}{\mathrm{HOT}^{\mathrm{lc}}_W}
\newcommand{\Hot}{\mathrm{Hot}}
\newcommand{\Hom}{\mathrm{Hom}}
\newcommand{\hot}{\mathrm{hot}}
\newcommand{\id}{\mathrm{id}}
\newcommand{\Wd}{W^{\mathrm{d}}_S}
\newcommand{\Wf}{W^{\mathrm{f}}_S}

\begin{document}

\printedpage{34}
% head: A. Grothendieck : Dérivateurs, Ch. XII, Édition des Universités de Montpellier II et Paris VII
% foot: Version 7 décembre 2002

\msp{page 46}
so if $f,g : S' \rightrightarrows S$ are two homotopic functors, then
\[ f^* \simeq g^* : \HOTlc(S) \rightrightarrows \HOTlc(S'). \]

The following follows at once

\textbf{Corollary 2.} \textit{If $f$ is a homotopism, $f^* : \HOTlc(S) \longrightarrow \HOTlc(S')$ is an equivalence of categories.}

Indeed, taking $g : S \longrightarrow S'$ such that
\[ fg \sim \id_S, \qquad gf \sim \id_{S'}, \]
we find that
\[ \begin{aligned}
g^*f^* \simeq (\id_S)^* &= \id_{\HOTlc(S)} \\
f^*g^* \simeq (\id_{S'})^* &= \id_{\HOTlc(S')},
\end{aligned} \]
so we find that (adjusting these functorial isomorphisms if need be) $f^*$ and $g^*$ are equivalences quasi-inverse to one another.

I have the impression (\footnote[54]{O.K.}) that all this should not use more than the standard axioms on $W$ as a basic localizer, at most W(2) instead of the weak form W(2 bis).

Let us return to $X$ smooth over contractible $S$. Corollary 2 implies that
\[ \begin{aligned}
\Hot_W &\longrightarrow \HOTlc(S) \\
\xi &\longmapsto p^*_S(\xi),
\end{aligned} \]
where $p_S : S \longrightarrow e$, is an equivalence. So in the category $\HOTlc(S)$, $X$ is isomorphic \msp{page 47} to a Cat-fibered object of the form $Z \times S \xrightarrow{\mathrm{pr}_2} S$, with $Z$ in Cat. Now for $X_1 = Z \times S \longrightarrow S$, it is immediate that $X_{1\,s} \longrightarrow X_1$ is in $W$, for a Cartesian product of arrows \texttt{[in]} $W$ (namely $\id_Z$ and $s \longrightarrow S$) is in $W$. (In fact, the inclusion is even a homotopism!) Let us now note that if $X_0$, $X_1$ are two objects in $\underline{\mathrm{Liss}}_0S$, such that there exists an $S$-arrow $X_0 \longrightarrow X_1$ which is in $\Wd$, then $X_{0\,s} \longrightarrow X_0$ is in $W$ if and only if $X_{1\,s} \longrightarrow X_1$ is. Indeed, we have
\[ \begin{tikzcd}
X_{0\,s} \arrow[r, "\in W"] \arrow[d, "i_s^{X_0}"'] & X_{1\,s} \arrow[d, "i_s^{X_1}"] & \\
X_0 \arrow[r, "\in W", ""{name=A, below}] & X_1 & \\
 & & \text{\footnotesize\begin{tabular}{@{}l@{}}here we probably use\\ W(2), not only W(2 bis),\end{tabular}}
\arrow[from=A, to=3-3, bend right]
\end{tikzcd} \]
\printedpage{35}
% head: A. Grothendieck : Dérivateurs, Ch. XII, Édition des Universités de Montpellier II et Paris VII
% foot: Version 7 décembre 2002
which proves the result.

It is true that the fact that $X = X_0$ and $X_1 = Z \times S$ are isomorphic in $\HOTlc(S)$ does not imply, by itself, the existence of a \msp{page 48} chain of isomorphisms coming from arrows in $\Wd$. But in fact, if in the statement of corollary 2 above we compare $X$ in $\underline{\mathrm{Fib}}_0S$ with $g^*f^*(X)$, thanks to a homotopy $\id_S \sim fg$, we do indeed find a sequence of functors $\varphi_i : S \longrightarrow S$, with $\varphi_0 = \id_S$, $\varphi_N = fg$, and $\varphi_i$ and $\varphi_{i+1}$ linked by an arrow in $\underline{\mathrm{Hom}}(S,S)$, so we have a sequence of objects $X_i = \varphi_i^*(X)$ of $\underline{\mathrm{Fib}}_0S$ (with $X_0 = X$, $X_N = g^*f^*X$), and $X_i$ and $X_{i+1}$ linked by an arrow \textit{in} $\underline{\mathrm{Fib}}_0S$ (not only in $\HOTlc(S)$), which arrow is in $\Wd = \Wf$. So we win, subject to various routine checks (\footnote[55]{O.K.}). So modulo these, we will have

\textbf{Proposition 9} (\footnote[56]{It's ok assuming that $W$ is \textit{strongly} saturated, i.e. $\hot_W(f)$ isomorphism $\Longrightarrow f \in W$.},\footnote[57]{So left $W$-fibration $\Longleftrightarrow$ right $W$-fibration $\Longleftrightarrow$ $W$-fibration.}). \textit{Suppose that $X$ in $\mathrm{Cat}/S$ is a left $W$-fibration (cf. prop. 8). Then it is even a $W$-fibration (in the sense of b), page 40, where it is called ``$W$-fibration \textup{hot}-strong''), in other words, if $S'' \longrightarrow S'$ is over $S$ with $S''$, $S'$ contractible, then $X'' \longrightarrow X'$ is in $W$ (\footnote[58]{we only find that $\hot_W(X'') \longrightarrow \hot_W(X')$ \texttt{[is an]} isomorphism.}). We probably have better: if $S'' \longrightarrow S'$ is a homotopism, then $X'' \longrightarrow X'$ is in $W$.}

\msp{page 49}
Let us prove (if we can) this nice complement. I point out that I had to make the hypothesis of $W$-presmoothness (or of $W$-preproperness) to be sure that the inverse image $X \times_S S'$, under $f' : S' \longrightarrow S$, really does coincide with $f^*$ on the $\HOTlc$ - otherwise I have no chance. % [sic] où -> ou, the conjunction "or" is meant
``We know'' (everything has to be looked at again!) that it's O.K. in the three cases $X$ $W$-smooth, $X$ $W$-proper, $X$ a $W$-fibration (but in what sense???). The check to be made (when working with the $\Wd$) is this:
\[ \begin{aligned}
&X \xrightarrow{i_X} \Phi_S(X) = Y \text{ is in } \Wd \text{ (without merit), and} \\
&\text{remains so under any base change.}
\end{aligned} \tag*{$(*)$} \]

But for $X, Y \in \mathrm{Ob}\,\mathrm{Cat}/S$ which are $W$-presmooth, $\Wd$-equivalence coincides with $\Wf$, which is preserved by base change, as is the hypothesis on $X$, $Y$. So we can rest easy.

We are reduced to proving that if $X$ is $W$-presmooth over $S$ and a $W$-pseudofibration, and if $S' \longrightarrow S$ is a homotopism, then $X' \longrightarrow X$ is in $W$. By the earlier argument using $X \longrightarrow \Phi_S(X)$, we reduce to the case where $X$ is split fibered.

\printedpage{36}
% head: A. Grothendieck : Dérivateurs, Ch. XII, Édition des Universités de Montpellier II et Paris VII
% foot: Version 7 décembre 2002

\msp{page 50}
\[ \begin{tikzcd}
X' \arrow[rr] \arrow[d] & & X \arrow[d] \\
S' \arrow[rr, "f"] \arrow[dr, "p'"'] & & S \arrow[dl, "p"] \\
 & e &
\end{tikzcd} \]

The arrow $X' = f^*(X) \longrightarrow X$ is interpreted in $\Hot_W$ as
\[ p'_!(f^*(\xi)) \longrightarrow p_!(\xi) \ , \qquad \text{where } \xi = \hot_{S,W}(X). \tag*{$(*)$} \]

I say that this arrow is \texttt{[an]} isomorphism, i.e. that for every $\eta \in \Hot_W$, the transposed arrow
\[ \Hom_W(p_!(\xi),\eta) \longrightarrow \Hom_W(p'_!(f^*(\xi)),\eta) \tag*{$(**)$} \]
is \texttt{[a]} bijection. Granting the adjunction formula, the arrow $(**)$ is interpreted as
\[ \Hom_W(\xi,\underbrace{p^*(\eta)}_{\eta_S}) \longrightarrow \Hom_W(f^*(\xi),\underbrace{p'^*(\eta)}_{f^*(\eta_S)}), \]
i.e.
\[ \Hom_W(\xi,\eta_S) \longrightarrow \Hom_W(f^*(\xi),f^*(\eta_S)), \]
that this is \texttt{[an]} isomorphism therefore follows from $f^*$ fully faithful on the $\HOTlc$. Now we have seen that $f^*$ is an equivalence (cor. 2, page 46).

\msp{page 51}
I am going to take up propositions 9 and 8 again, a bit badly put together.

\textbf{Theorem 3.} \textit{We assume that $W$ satisfies \textup{W(1,2,3)}. Let $f : X \longrightarrow S$. Consider the following conditions.}

\begin{itemize}[leftmargin=4em]
\item[\textit{a)}] \textit{$\forall\ S'' \longrightarrow S'$ over $S$, with $S''$, $S'$} contractible, \textit{$X'' \longrightarrow X'$ is in $W$.}

\item[\textit{a}$'$\textit{)}] \textit{$\forall\ S'$ over $S$, with $S'$ contractible, and $s' \in S'$, the inclusion} % AUDIT resolved: the label is a prime, not d prime; see corrections.csv, PDF leaf 46
\[ X_{s'} \hookrightarrow X' = X \times_S S' \]
\textit{is in $W$.}

\item[\textit{a}$_1$\textit{)}] \textit{As in a$'$), but assuming that $S'$ has a final object (instead of $S'$ contractible).} % AUDIT resolved: the reference is a prime; see corrections.csv, PDF leaf 46

\item[\textit{a}$_2$\textit{)}] \textit{As in a$'$), but assuming that $S'$ has an initial object.} % AUDIT resolved: the reference is a prime; see corrections.csv, PDF leaf 46

\printedpage{37}
% head: A. Grothendieck : Dérivateurs, Ch. XII, Édition des Universités de Montpellier II et Paris VII
% foot: Version 7 décembre 2002

\item[\textit{b}$_1$\textit{)}] \textit{$f$ is $W$-}preproper \textit{(i.e. a$'$) is satisfied when $s'$ is a final object of $S'$, with no other condition on $S'$ in $\mathrm{Cat}/S$), and moreover for every arrow $s_0 \longrightarrow s_1$ in $S$, the induced arrow $X/s_0 \longrightarrow X/s_1$ is in $W$ (i.e. a$'$ \textup{\texttt{[rather a]}}) is satisfied for $Y/s_0 \longrightarrow Y/s_1$ \textup{\texttt{[rather} $S/s_0 \longrightarrow S/s_1$\texttt{]}}).}

\item[\textit{b}$_2$\textit{)}] \textit{$f$ is $W$-}presmooth \textit{(i.e. a$'$) is satisfied when $s'$ is an initial object of $S'$, with no other condition on $S'$), and moreover the arrows $s_0\backslash X \longrightarrow s_1\backslash X$ \textup{\texttt{[rather} $s_1\backslash X \longrightarrow s_0\backslash X$\texttt{]}} (for $s_0 \longrightarrow s_1$ in $S$) are in $W$, (i.e. a$'$ \textup{\texttt{[rather a]}}) satisfied for $s_1\backslash Y \longrightarrow s_0\backslash Y$ \textup{\texttt{[rather} $s_1\backslash S \longrightarrow s_0\backslash S$\texttt{]}}).}

\item[\textit{c}$_1$\textit{)}] \textit{$f$ is $W$-preproper and a $W$-pseudofibration, i.e. condition a$'$) is satisfied in the following two cases: 1$^\circ$) $s'$ is a final object of $S'$, 2$^\circ$) $S' = \mathbf{\Delta}^1$.}

\item[\textit{c}$_2$\textit{)}] \textit{$f$ is $W$-presmooth and a $W$-pseudofibration, i.e. a$'$) in the two cases: 1$^\circ$) $s'$ is an initial object of $S'$, 2$^\circ$) $S' = \mathbf{\Delta}^1$.}
\end{itemize}

\textit{All these 8 conditions are equivalent. They imply the following:}

\begin{itemize}[leftmargin=4em]
\item[\textit{d)}] \textit{Let $S'' \xrightarrow{u} S'$ be a homotopism in $\mathrm{Cat}/S$, and consider $u_X : X'' \longrightarrow X'$. Then $\hot_W(u_X)$ is \textup{\texttt{[an]}} isomorphism. (So if $W$ is strongly saturated, $u_X \in W$. So under this condition on $W$, d) is equivalent to the preceding conditions.)}
\end{itemize}

\msp{page 52}
\textbf{Definition 4.} We say that $f : X \longrightarrow Y$ in Cat is a \textit{$W$-fibration}, if it satisfies the equivalent conditions a) to c$_2$) of proposition 9.

\textbf{Corollary 1.} \textit{If $f$ is $W$-proper or $W$-smooth, for $f$ to be a $W$-fibration, it is necessary and sufficient that $f$ be a $W$-pseudofibration, i.e. that for $S' = \mathbf{\Delta}^1 \longrightarrow S$, the inclusions of the two fibers}
\[ X'_0 \xrightarrow{i_0} X' \xleftarrow{i_1} X'_1 \]
\textit{of $X' = X \times_S S'$ be in $W$.}

\textbf{Corollary 2.} \textit{Suppose $f$ \textup{Cat}-fibrant. Then $f$ is a $W$-fibration if and only if for every arrow $u : s \longrightarrow s'$ in $S$, $u^* : X_{s'} \longrightarrow X_s$ is in $W$. Dual statement if $f$ is \textup{Cat}-cofibrant.} % [sic] pout -> pour, misprint

\textbf{Corollary 3.} \textit{Suppose $f$ $W$-preproper and $W$-presmooth (for example $f$ $W$-perfect). Then $f$ is a $W$-fibration.}

\textbf{Proposition 10.} \textit{Let $f : X \longrightarrow Y$ be an arrow in $\underline{\mathrm{Fib}}_WS$ (\footnote[59]{$X$ and $Y$ both preproper or both presmooth suffices.}) (full subcategory of $\mathrm{Cat}/S$ formed by the $W$-fibered objects over $S$). If $\forall\ s \in S$, $f_s : X_s \longrightarrow Y_s$ is in $W$, then $f \in W$.}

Let us consider
\printedpage{38}
% head: A. Grothendieck : Dérivateurs, Ch. XII, Édition des Universités de Montpellier II et Paris VII
% foot: Version 7 décembre 2002
\[ \begin{tikzcd}
X \arrow[rr, "f"] \arrow[dr, "i_X"'] & & Y \arrow[dr, "i_Y"] & \\
 & \Phi_S(X) \arrow[rr, "\Phi_S(f)"'] & & \Phi_S(Y),
\end{tikzcd} \]
where $i_X, i_Y \in W$, so $\Phi_S(f) \in W \Rightarrow f \in W$. Now the fact that $X$ \texttt{[is]} $W$-fibered over $S$ implies that $i_X \in \Wf$ (it even suffices that $X$ be preproper \texttt{[rather presmooth]} over $S$), likewise for $Y$, $i_Y$. So the hypothesis $f \in \Wf$ implies $\Phi_S(f) \in \Wf$, so $\Phi_S(f) \in W$ by \texttt{[the]} proposition 5 (page 29).

\end{document}
